\chapter{Introduction}
Over the last century, phones have changed from big and clumsy boxes that were mostly used to transfer simple spoken information over long distances into something different. Nowadays, the capabilities of smartphones allow us to navigate through the streets of big cities we have never been to before, translate conversations between people speaking different languages,
% edited(
and recognize objects, providing us, the users, with information about them in a matter of seconds.

Board games have likewise changed from the simplest ones, such as checkers, which can be played on any surface with a checkerboard pattern and rocks of different colors or shapes, into rather complex games. Applications helping players track their statistics for board games have existed for some time. A great example would be \textit{D\&D Beyond}, which allows players to keep and manage their in-game characters for \textit{Dungeons and Dragons}. Players would otherwise have to rely on a 380 page rule-set book dedicated to character information and game functionalities. 

Further advancements in this field, especially for smartphones advancing in performance each year, seem inevitable. An application utilizing computer vision for the game \textit{Krycí Jména}, was developed back in 2018 \cite{hurta} by one of my predecessors. Applications combining computer vision and board games ease the learning curve of games and shorten the time it takes for new players to grasp the basics of the game. These applications were referred to as smart assistants by one of my predecessors \cite{takacs}, who also explored this topic.

To further explore this topic, an end-to-end object detection pipeline was proposed. Its first part, Computer Vision Inference Library for Board Games, or CVLiBG, is an Android library written in Kotlin, whose purpose is to provide a base for developers to build their applications upon. The aim is to provide tools for creating applications utilizing computer vision to detect board game objects. Furthermore, the library contains template functions for loading assets, storing user data, and other useful tools for repetitive tasks.

Object detection models for common objects, such as fruits or everyday objects, are usually widely available on the internet. Game specific models, on the other hand, might not always be available, and neither may the datasets containing these game objects. Because of that, a desktop application called Dataset Creator was created. This tool allows developers to take raw image data, annotate it, store it, increase the size of the dataset by using augmentation, and finally train it within the application. Trained models can be further used for annotating additional image data.

The proposed solution should help developers create annotated datasets from raw image data obtained by the developers. Furthermore, the mobile library lets developers focus on graphical design and interactions with users, while object detection and camera handling are addressed by the library. The resulting applications are capable of recognizing the game objects on which the model was trained. This thesis also demonstrates the use of the proposed tools to create an application for the Bang! card game. The final application is capable of showing descriptions with tips or real life examples based on recognized cards. All of this is achieved with a custom object detection model trained on a dataset created within the proposed Dataset Creator tool. While the image data was obtained by taking photos of the card game.

The object detection, as a method of computer vision, is very capable and is usually the go-to method for real-time object detections. However, this method struggles with complicated scenes containing multiple, especially small, objects. The proposed solution solves camera control and object detection by making development faster and more accessible to hobbyists and small developers, but it does not replace the overall process of creating these applications.

The next chapter explores the general knowledge about computer vision, object detection, neural networks, and how to train them. Followed by the chapter explaining the techniques used in object detection. The fourth chapter describes the implementation and design of the CVLiBG process. The next chapter describes the implementation of the Dataset Creator and the data augmentation techniques it provides. (\textbf{TBD: add experiments})


% edited )